\documentclass[11pt,a4paper, top=1cm]{article}
\usepackage[a4paper,left= 2cm, right= 2cm, top= 1.5cm, bottom= 1cm, includefoot]{geometry}

\usepackage[T1]{fontenc}
\usepackage[utf8]{inputenc}
\usepackage[ngerman]{babel}
\usepackage{amsmath}
\usepackage{enumitem}
\usepackage{parskip}
\usepackage{titlesec} \titleformat{\section}{\normalfont\large\bfseries}{\thesection}{1em}{}
\titleformat{\subsection}{\normalfont\normalsize\bfseries}{\thesubsection}{1em}{}
\titlespacing{\section}{0pt}{1.5ex plus .2ex}{1ex}
\setlist{nosep,leftmargin=1.4em}
\usepackage[unicode,hidelinks]{hyperref}
\hypersetup{
    pdftitle={Othello-Spieler – Kurzbeschreibung der Lösung},
    pdfauthor={Dominik Wendland},
    pdfsubject={Ausarbeitung},
    pdfkeywords={Othello, Minimax, NegaScout},
    pdflang={de}
}

\title{\vspace{-2.4em}Othello-Spieler – Kurzbeschreibung der Lösung\\
\vspace{0.2cm} \large Modul: Intelligente Systeme
\\\vspace{0.2cm} \small Betreuer: Prof. Dr. J. Priemer\\
\author{\large Dominik Wendland (\texttt{spieler.wendo})}\\
\date{\today}}

\begin{document}
\maketitle
\vspace{-2em}

\section{Überblick}
Der Spieler implementiert die vorgegebene Schnittstelle \texttt{OthelloSpieler}.
Die Zugberechnung erfolgt mit Hilfe eines Spielbaums durch eine
\textit{Minimax-Suche}\footnote{S. J. Russell und P. Norvig, Künstliche Intelligenz: ein moderner Ansatz, 4. Aktualisierte Auflage, Pearson, 2023.},
die in der \textit{Negamax}-Formulierung umgesetzt ist.
Das \textit{Alpha-Beta-Pruning} ist dabei als
\textit{NegaScout}\footnote{A. Reinefeld, Spielbaum-Suchverfahren, Bd. 200. in Informatik-Fachberichte, Springer Berlin Heidelberg, 1989} realisiert.
Als \textit{iterative Tiefensuche} unter Zeitkontrolle vertieft sich die Suche
schrittweise (Tiefe $1, 2, 3, \dots$), bis das Zeitbudget des Zuges erschöpft ist,
und liefert den besten Zug der letzten vollständig durchsuchten Tiefe.
Zur effizienten Beschneidung des Spielbaums werden die Züge an jedem Knoten heuristisch
sortiert -- an
jedem inneren Knoten nach statischer Feldbewertung und \textit{Killerzügen}, ergänzt um die
Prinzipalvariante aus der iterativen Tiefensuche.
Die Positionsbewertung ist phasenabhängig und gewichtet im Mittel- und Endspiel
unterschiedliche Merkmale.
Die folgenden Abschnitte gehen auf die einzelnen Aspekte der Lösung ein.


\section{Architektur}
Der Code ist nach dem Single-Responsibility-Prinzip modular in fünf Klassen mit
folgenden Aufgaben gegliedert:
\begin{itemize}
\item \texttt{Spieler} -- Adapter zur Wettkampf-Umgebung, übersetzt \texttt{Farbe} und
\texttt{Zug} in die interne Darstellung und zurück, pflegt das eigene Spielbrett und
behandelt Sonderfälle (Passen, ungültiger Gegnerzug $\to$ \texttt{ZugException}).
\item \texttt{Board} -- Spielregeln und Spielstand (Zuggenerierung, Ausführen inkl.
Umdrehen, Endstellungserkennung).
\item \texttt{MiniMax} -- Suchverfahren (NegaScout, Zugsortierung, iterative Tiefensuche,
Zeitabbruch).
\item \texttt{Evaluate} -- Positionsbewertung.
\item \texttt{TimeManager} -- Aufteilung der Gesamtbedenkzeit auf die einzelnen Züge.
\end{itemize}


\section{Brettdarstellung}
Das $8\times8$-Brett ist ein flaches \texttt{int}-Array der Länge~64 und
zeilenweise mit $\text{Index}=\text{Zeile}\cdot 8+\text{Spalte}$ indiziert
(Zeile~0, Spalte~0 ist das Feld oben links).
Jede Zelle trägt einen der Werte \texttt{EMPTY}, \texttt{BLACK} oder \texttt{WHITE} und
bildet so die Framework-Farben \texttt{Farbe.SCHWARZ} und \texttt{Farbe.WEISS} ab.
Die Suche folgt dem \textit{copy-make}-Prinzip.
Für jeden zu prüfenden Zug wird über \texttt{copy()} eine Kopie des Boards angelegt
und der Zug nur auf dieser Kopie ausgeführt, sodass die Ausgangsstellung unberührt bleibt.
Eine solche Kopie ist ein Array aus 64 Werten und wird als Ganzes per
\texttt{System.arraycopy} dupliziert.
Da die Alpha-Beta-Suche je Zug sehr viele Knoten expandiert und an jedem Knoten eine Kopie
erzeugt, muss das Duplizieren einer Stellung billig sein, was ein flaches Primitiv-Array leistet.


\section{Zugberechnung}
Bei jedem Aufruf von \texttt{berechneZug} wird zunächst der letzte Zug des Gegners auf dem
eigenen Brett synchronisiert (\texttt{replayOpponentMove()}), anschließend
über \texttt{TimeManager} das Zeitbudget für diesen Zug bestimmt und die Suche gestartet.

\paragraph{Suchverfahren (NegaScout).}
Die Suche ist als Negamax formuliert.
Beide Spieler maximieren aus eigener Sicht, und bei
jedem Halbzug wird das Vorzeichen der Kindbewertung negiert.
Darauf setzt NegaScout auf.
Der laut Zugsortierung vielversprechendste Zug wird mit vollem Suchfenster untersucht, jeder
weitere zunächst nur mit einem Nullfenster als billige Probe „ist dieser Zug überhaupt
besser?“.
Erst wenn die Probe dies andeutet (\textit{fail high}), wird der Zug mit vollem
Fenster nachgesucht.
Bei guter Zugsortierung expandiert NegaScout dadurch deutlich weniger
Knoten als einfaches Alpha-Beta, denn die meisten Geschwisterzüge werden schon vom billigen
Nullfenster widerlegt und nur die wenigen tatsächlich besseren Züge kosten eine volle
Nachsuche.

\paragraph{Zugsortierung.}
Alpha-Beta bzw.\ NegaScout beschneiden den Spielbaum nur dann effizient, wenn der beste Zug
zuerst untersucht wird.
An jedem inneren Knoten wird die Zugreihenfolge daher heuristisch bestimmt:
\begin{itemize}
\item \textbf{Statische Feldbewertung} -- Vorsortierung nach einer festen Gewichtstabelle
(Ecken sehr hoch, ecknahe „X“/„C“-Felder negativ).
\item \textbf{Killerzüge} -- pro Suchtiefe werden bis zu zwei Züge gemerkt, die
zuletzt einen Beta-Schnitt ausgelöst haben, und an Geschwisterknoten zuerst probiert.
Ein Zug, der einen Geschwisterknoten widerlegt, widerlegt häufig auch die übrigen.
Auf einer festen Stellung mit exakter Knotenzählung (\texttt{NodeBenchmark}) senkte diese
Maßnahme die Knotenzahl im Mittelspiel bei Tiefe~10 um
rund 37\,\% ($2{,}27$ auf $1{,}43$ Mio.).
\end{itemize}
Über die Iterationen hinweg tritt an der Wurzel die Prinzipalvariante hinzu: Der
beste Zug der vorherigen, flacheren Suche wird in der nächsten Iteration zuerst probiert.

\paragraph{Iterative Tiefensuche und Zeitmanagement.}
Statt sofort mit fester Tiefe zu suchen, wird nacheinander mit Tiefe $1,2,3,\dots$ gesucht.
Das hat zwei Vorteile: Es liegt jederzeit ein vollständiges Ergebnis vor (wichtig für die
Zeitkontrolle), und jede flachere Iteration verbessert die Zugsortierung der nächsten.
\texttt{TimeManager} weist jedem Zug ein Budget aus der verbleibenden Bedenkzeit geteilt
durch die geschätzte Restzugzahl (etwa die Hälfte der freien Felder) zu, vermindert um eine
Sicherheitsreserve von 10\,\%. Läuft das Budget während einer Iteration ab, bricht die Suche
über eine \texttt{TimeUpException} sofort ab und liefert den besten Zug der letzten
vollständig abgeschlossenen Tiefe. Die Tiefe ist dabei nur eine Obergrenze, real begrenzt die
Zeit. Im Endspiel mit wenigen freien Feldern wird die Partie faktisch bis zum Ende
ausgerechnet.


\section{Positionsbewertung}
Die Bewertungsfunktion \texttt{evaluate()} liefert einen Wert aus Sicht der ziehenden
Seite und ist antisymmetrisch
($\text{bewerte}(b,\text{seite}) = -\text{bewerte}(b,\text{gegner})$), was die
Negamax-Formulierung der Suche voraussetzt.
Sie ist phasenabhängig und wechselt bei 50 Steinen auf dem Brett von einer heuristischen
Mischung zu einer zählenden Bewertung.
\begin{itemize}
\item \textbf{Mittelspiel} (unter 50 Steinen): gewichtete Summe aus \textit{Eckenkontrolle},
\textit{Mobilität} (normierte Differenz der Zugmöglichkeiten) und \textit{Positionswert} (Feldtabelle).
Die Ecken tragen das größte Gewicht, denn
ein besetzter Eckstein kann nie wieder umgedreht werden und ist damit dauerhaft sicher.
Die reine Steinzahl bleibt bewusst unberücksichtigt, weil ein früher Steinvorsprung in
Othello meist nachteilig ist, da er die eigene Mobilität einengt und dem Gegner
Möglichkeiten zum Umdrehen liefert.
\item \textbf{Endspiel} (ab 50 Steinen): Die Mischung wird durch die \textit{Steindifferenz}
ersetzt, weiterhin mit starker Gewichtung der Ecken. Nun entscheidet die endgültige
Steinzahl über Sieg oder Niederlage, sodass Mobilität und Feldtabelle ihren Vorhersagewert
verlieren und stattdessen direkt gezählt wird.
\end{itemize}


\section{Fazit}
Der Spieler verbindet eine NegaScout-Suche mit iterativer Tiefensuche, dreistufiger
Zugsortierung und einer phasenabhängigen Bewertung.
Leistungsänderungen wurden an exakten Knotenzahlen
überprüft und nur bei nachgewiesenem Nutzen übernommen, das Verhalten ist durch
48 JUnit-Tests abgesichert.
Im Wettbewerb belegte der Spieler dennoch nur den vierten Platz -- weniger, als die
Optimierung der Suche erwarten ließ.
Die folgende Einordnung ist eine Hypothese, keine Diagnose.
Die Suchoptimierung gewann über die Geschwindigkeit zwar mehr Tiefe, doch dürfte die von
Hand gewichtete Bewertung nicht ausreichend abgestimmt gewesen sein.
Die vielversprechendsten Hebel wären entsprechend beispielsweise eine sorgfältiger abgestimmte
oder gelernte Bewertung und ein früher exakt gelöstes Endspiel.
Welcher davon tatsächlich entscheidet, könnte erst eine Analyse der Turnierpartien zeigen.

\end{document}
